Language models that are told a secret and asked not to reveal it still leak it into the stories they write.
Would an explicit plan relieve the leak?
We run a white-box study of \gemma: the model writes stories while holding a secret word, a guesser decides which of two stories was written with a given secret (chance is 50\%), and we read and edit the residual stream while the model writes.
A plan removes the leak only if its content was fixed without knowledge of the secret.
With no plan the guesser is right 84.1\% of the time.
An outline that the same model wrote without the secret brings this to 49.7\%, and a one- or two-sentence secret-blind premise to 54.7\%.
An outline the model writes while holding the secret leaks by itself (83.8\%) and passes the leak on to the story (81.9\%).
Token-matched context that fixes no content leaves leakage at 75.1\%.
Inside the model, the secret is decodable at story positions with 97 to 100\% accuracy (15-way) in every condition, including the one with no behavioral leak, so representing a secret is not the same as leaking it.
Removing the secret from the context after 30 story tokens abolishes the leak (52.2\%), while withholding it for the first 100 tokens does not (77.5\%): the leak is produced continuously as content is chosen.
Results for plot twists point the same way with weaker evidence.
To keep known information out of generated prose, decide the content before the model sees the secret.
